\section{From 3D uplifting to downstream tasks}

In this section, we describe our pipeline around the uplifting process from Sec.~\ref{sec:method_uplifting}, covering preceding and following stages for three tasks: multi-view segmentation, open-vocabulary object segmentation, and open-vocabulary semantic segmentation. 
As in Sec.~\ref{sec:method_uplifting}, we are given a set of 2D frames $I_1, \dots, I_m$, with viewing directions $d_1, \dots, d_m$, and a 3D Gaussian Splatting representation of the scene.

\subsection{Multi-view segmentation}\label{sec:muti_seg}
\label{sec:method_seg}
We assume that a foreground mask of the object to be segmented is provided on the \emph{reference frame}~$I_1$.
This mask can either be a full segmentation or a sparse set of \emph{scribbles}, both defining a set of foreground pixels $\mathcal{P}$.
The task is to generate a 2D segmentation mask on one or more \emph{target frames} based on the foreground from the \emph{reference frame}. It is evaluated with the intersection over union (IoU). 

\myparagraph{Segmentation with SAM}
SAM~\citep{kirillov2023sam, ravi2024sam2} is a supervised pretrained model that, given an input image and point prompts, generates segmentation masks. 
We uplift SAM masks from multiple views using Eq.~(\ref{eq:uplifting}); this aggregation improves cross-view consistency and leads to better single-view segmentation. On each view, SAM is run multiple times with different prompt sets obtained by uplifting and reprojecting the reference mask, and the resulting predictions are averaged (see Sec.~A.1).
The final prediction is obtained by rendering the 3D mask into the target frame and thresholding.
Results obtained with this strategy are reported for SAM~\citep{kirillov2023sam} and SAM2~\citep{ravi2024sam2} in Sec.~\ref{sec:exp-seg}. 

\myparagraph{Segmentation with DINOv2}
We construct 2D feature maps using DINOv2 with registers~\citep{darcet2024registers}, applying a sliding window over the image followed by dimensionality reduction of the DINOv2 features.
To favor the first principal components, known to focus on the foreground objects~\citep{oquab2024dinov2}, the features are re-weighted by the eigenvalues of the PCA decomposition.
The 2D feature maps from the $m$ training views are uplifted using Eq.~(\ref{eq:uplifting}) and the resulting 3D features are re-projected into any direction using Eq.~(\ref{eq:rendering}).
2D segmentation is obtained by thresholding the similarity between projected features and those from the foreground pixels $\mathcal{P}$ in the reference frame (see Sec.~A.2).
This approach corresponds to an ablation discussed in Sec.~\ref{sec:exp-seg}.
%

\myparagraph{Improving DINOv2 segmentation with graph diffusion}
Segmentation based on DINOv2 alone may struggle to distinguish objects with similar features, as it lacks 3D spatial information.
To address this, we leverage the graph diffusion process introduced in Sec.~\ref{sec:diffusion} and illustrated in Fig.~\ref{fig:diffusion}. 
The foreground mask from the reference frame is first uplifted into 3D, defining a set of anchor Gaussians $\mathcal{M}$. This set is used to initialize the weight vector $g_0$ and to define the regularization term $P$, based on the similarity between each 3D feature and those of Gaussians in $\mathcal{M}$ (see supplementary Sec.~\ref{sec:appendix-diffusion}).
For this task, we binarize $A$ with a fixed threshold (set to $10^{-5}$). After $T$ diffusion steps, we recover the  nodes $\mathcal{S}$ in $g_T$ with strictly positive values (\emph{i.e.}, those reachable after $T$ steps). The final weight is defined as $h_i = P(f_i)$ if $i\in \mathcal{S}$ and $0$ otherwise. Segmentation is then performed by projecting $\mathbf{h} = (h_i)$ into 2D and thresholding.
This approach achieves results comparable to SAM, as reported and discussed in Sec.~\ref{sec:exp-seg}.

\subsection{Open-vocabulary object segmentation}
\label{sec:method_clip}

Following \cite{kerr2023lerf}, we tackle the task of open-vocabulary object localization by uplifting CLIP features~\citep{ilharco2021openclip}.
CLIP effectively aligns images and text in a shared representation space. As a measure of alignment, we use the relevancy score introduced by LERF~\citep{kerr2023lerf}, which measures the similarity between a CLIP visual feature and a text query.

\myparagraph{Construction of CLIP feature maps}
We follow common practice \citep{kerr2023lerf, zuo2024fmgs} and construct multi-resolution CLIP 2D feature maps by querying CLIP on a grid of overlapping patches at different scales. 
As in \cite{zuo2024fmgs}, rather than keeping the resulting representations separate, we aggregate them via average pooling. These multi-resolution CLIP features are uplifted into 3D using Eq.~(\ref{eq:uplifting}).

\myparagraph{Relevancy scores} 
After uplifting CLIP features, we compute relevancy scores for each Gaussian's feature to text queries embedded by CLIP. These relevancy scores can then be projected into 2D and used for both localization and segmentation. For localization, we choose the pixel with the highest relevancy score.
For segmentation, we render the relevancy scores and either (i) directly threshold the rendered scores, or (ii) predict a SAM mask by selecting point prompts among pixels with the highest score. Specifics on the computation of relevancy scores and segmentation masks are provided in the supplementary (Sec.~\ref{sec:appendix-clip}).

\myparagraph{Refining relevancy scores with DINOv2 graph diffusion}
We refine 3D relevancy scores with the diffusion process described in Sec.~\ref{sec:diffusion}. To this end, DINOv2 features are also uplifted, and the similarity matrix is built as in Eq.~(\ref{eq:edges}), with the unary term $P$ constructed using a logistic regression over thresholded relevancies, see details in the supplementary (Sec.~\ref{sec:appendix-clip}). The diffusion process propagates CLIP relevancies to Gaussians with similar DINOv2 features. The resulting 3D relevancy scores span the object of interest without covering other objects with similar features and show a strong decay at the object's borders, as defined by DINOv2's feature landscape, resulting in improved segmentation results. 
This approach is evaluated and compared to direct segmentation without diffusion in Sec.~\ref{sec:exp-clip}.

\subsection{Open-vocabulary semantic segmentation}
\label{sec:method_semantic}

In semantic segmentation, one needs to label each Gaussian with the appropriate text label from a predefined set.
%
For this task, we directly replace OpenGaussian’s 3D feature learning with our uplifting approach, keeping the rest of their protocol unchanged. This enables a fair comparison under the same evaluation setup.
In particular, 2D feature map generation is performed using SAM in automatic mask generation mode and assigning a CLIP feature to each mask, as in~\citep{qin2023langsplat, shi2024legaussians, wu2024opengaussian}.
After uplifting, each Gaussian is assigned 
the text label with the highest CLIP similarity.
Results are reported in Sec.~\ref{sec:exp-semantic}.